\chapter{Background}

This chapter introduces the main background technologies used in the proposed blockchain-based reporting framework. It explains the basic concepts of blockchain technology, permissioned networks, PoA consensus, smart contracts, decentralized applications, \ac{IPFS} , privacy-preserving identity concepts, and AI-based content moderation. Through this chapter we aim to give the reader enough technical understanding to follow the system design and implementation discussed in the later chapters.

\section{Blockchain Technology}


Blockchain is a distributed, digital, append-only ledger technology. It enables data to be stored in immutable and transparent manner across multiple nodes in a network. Unlike centralized systems, blockchain systems maintain a synchronized copy of the ledger among all authorized participants. This approach reduces the dependency on a single intermediary controlling authority. Each block in the blockchain contains a set of validated transactions, a timestamp, and a cryptographic hash of the previous block. This hash-chaining mechanism ensures that any alteration to previous recorded data becomes immediately detectable.

Immutability is a key characteristic in blockchain technology. It is computationally impractical to alter or delete a transaction after it is confirmed and appended to the blockchain. In systems that deal with governance, the immutable property is crucial for ensuring integrity and transparency.

Another key component of blockchain systems is Cryptography. Hash functions are used to make digital fingerprints that are unique for each piece of data. And also digital signatures are used to verify the authenticity and the ownership of the transactions. Therefore the use of cryptographic methods in blockchain systems guarantee the integrity, authenticity and non-repudiation of data.


% Cryptography plays a significant role in blockchain systems. To generate unique digital fingerprints of data, hash functions are used. To verify ownership and authenticity of transactions, digital signatures are used. By using above cryptographic methods, blockchain systems guarantee integrity, authenticity and non-repudiation of data. 

\section{Permissioned Blockchain and PoA Consensus}

% Bitcoin and Ethereum are permissionless blockchains, where anyone can participate as a node and validate transactions. In contrast, permissioned blockchains restrict participation to a predefined set of entities. These entities are often referred to as validator nodes. They are responsible for validating transactions and maintaining the integrity of the blockchain. 
%
% While permissionless blockchains rely on consensus mechanisms like \ac{PoW} or \ac{PoS}, the \ac{PoA} consensus mechanism is specially designed for permissioned blockchains. In here instead of relying on computationally expensive mining or staking process the validator nodes in the permissioned blockchain are validating blocks. Therefore this approach is a reputation and indetity-based mechanism, where the authority of the validators are based on their identity and reputation within the network. This allows for faster transaction processing and lower energy consumption compared to traditional consensus mechanisms. This makes permissioned blockchains more suitable for enterprise applications and private networks.

% Bitcoin and Ethereum are public, permissionless networks. A participant can run a node and help validate transactions without first obtaining approval from a central administrator. Permissioned blockchains follow a stricter model. Only a predefined set of organisations may join, and these members typically operate as validator nodes that verify transactions and append blocks to the shared ledger.

% Public chains usually reach agreement through \ac{PoW} or \ac{PoS}, where block production is tied to mining effort or staked cryptocurrency. \ac{PoA} was introduced for permissioned environments where every validator is already known to the network operator. Instead of competing for blocks, approved validators take turns signing them. Trust rests on verified identity and organisational accountability rather than hash power or token deposits. Confirmation times are shorter and the protocol avoids the energy overhead of mining, which is why many private and consortium deployments adopt permissioned \ac{PoA} rather than a public chain.

% Local governance related system requires predictable performance, immutability, transparency and accountability. \ac{PoA} consensus running in a permissioned blockchain is well suited for above requirements because validators identities are known and they controlled by trusted entities. In addition to that, \ac{PoA} eliminates the need for economic incentives. This make it more efficient and practical for local governance applications. 
Bitcoin and also Ethereum are examples of public, open blockchains. Anyone who wants to do so can become a participant and help out by running a node and verifying transactions without needing to first get approval from a central person or group. However, a more restricted model is used in permissioned blockchains. Only specific, approved organizations are allowed to join, and they generally act as special type of node that checks on and adds new groups of transactions to the shared record book.

Most public chains use either \ac{PoW} or \ac{PoS} consensus mechanisms. It means that making new blocks is based on how much processing power you have or how much of their own type of online money you're willing to use as a kind of stake. However, the \ac{PoA} was developed for use in more exclusive situations where all of the people or organizations that will be doing the verifying are already known to the person or organization in charge of the system. Unlike with \ac{PoW} and \ac{PoS}, these approved validators don't compete with each other by working to solve a problem, instead, they just take turns adding new blocks.This allows for faster transaction processing and lower energy consumption compared to traditional consensus mechanisms. This makes permissioned blockchains more suitable for enterprise applications and private networks.

Local governance related system requires predictable performance, immutability, transparency and accountability. \ac{PoA} consensus running in a permissioned blockchain is well suited for above requirements because validators identities are known and they controlled by trusted entities. In addition to that, \ac{PoA} eliminates the need for economic incentives. This make it more efficient and practical for local governance applications. 


\section{Smart Contracts}

Smart contracts are first introduced by Ethereum blockchain. They are self-executing contracts with the terms of the agreement directly written into code. In a blockchain based governance systems smart contract can be used to eliminate centralized administrative control. It can automate workflows that would traditionally require human intervention. Therefore this reduce the reliance on intermediaries, hence improve the transparency and trustworthiness.

The proposed system uses Solidity based smart contracts deployed on a private permissioned PoA blockchain.

\section{Decentralized Applications (\ac{DApp})}

\ac{DApp}s are software applications that utilize a decentralized network, such as a blockchain. Traditional web applications rely on centralized databases and backend infrastructure but \ac{DApp}s utilize smart contracts to handle backend logic and state management. This shift from Web2 to Web3 architecture ensures that application data is immutable, transparent, and resistant to single point of failure or centralized control.

% In the \ac{DApp} architecture, the frontend provides the user interface, while the core logic is distributed across the blockchain network. To interact with the blockchain, the frontend utilizes Web3 provider libraries, such as Web3.js or Ethers.js. These libraries enable the application to communicate with blockchain nodes via \ac{RPC}. Through this connection, the frontend can query the blockchain to read the current state of smart contracts or construct new transactions. When a user wishes to modify the blockchain state, they must cryptographically sign the transaction using their digital wallet's private key, ensuring authenticity and authorization.
%
% A significant challenge in traditional DApp adoption is the requirement for users to pay transaction fees, commonly known as ``gas," using the native cryptocurrency of the network. This creates a barrier to entry for regular users who may not possess or understand cryptocurrency. To address this usability issue, modern DApps employ meta-transactions and relayers. A meta-transaction allows a user to cryptographically sign their intent to execute a transaction off-chain without broadcasting it themselves. This signed message is then sent to an intermediary service known as a relayer. The relayer, which holds a balance of the network's cryptocurrency, wraps the user's signed message in a standard blockchain transaction, submits it to the network on behalf of the user, and pays the associated gas fees. This approach provides a seamless, gasless user experience while maintaining the security and non-repudiation guarantees of cryptographic signatures, making blockchain applications more accessible to the general public.

A typical \ac{DApp} splits work between a frontend and on-chain smart contracts. The frontend handles what the user sees and clicks through, while the rules and stored data sit on the blockchain. Developers link the two parts using Web3 provider libraries such as Web3.js or Ethers.js. Those libraries talk to blockchain nodes over \ac{RPC}, so the frontend can fetch the current state of a contract or build a new transaction. If the user wants to change on-chain data, the wallet must sign the transaction with its private key. Without that signature the network will not accept the transaction.

Paying transaction fees, or ``gas,'' in the network's own cryptocurrency is another problem for everyday users. Someone who does not already hold tokens or understand how wallets work may stop at the first fee prompt. MetaMask or similar wallet apps and backend relayers are one way around this. The user signs a message that describes the action they want, but does not broadcast it or pay gas themselves. Instead, the signed message goes to a relayer. A relayer is a service that keeps a funded account on the chain. The relayer packages the signature into a normal transaction, submits it, and pays the fee. The user still proves intent through cryptography, so the usual security properties hold, but the flow can feel more like a standard web application.


\section{\ac{IPFS}}

\ac{IPFS} is a peer-to-peer protocol for storing and sharing files across many nodes instead of a single central server. On the conventional web, a resource is found through a location based URL that points to a specific host. \ac{IPFS} identifies data by its content. Each file receives a \ac{CID}, which is a cryptographic hash derived from the file itself. If the content changes, the \ac{CID} changes too, so anyone retrieving the file can check that the bytes they received match the expected hash. Identical content always maps to the same \ac{CID}, which also allows deduplication across the network. Nodes may voluntarily store or pin data, and other peers can download a file from any node that still holds it. Availability therefore depends on at least one node continuing to pin that content.

Storing bulky files directly on a blockchain is expensive and slow, which is why many blockchain-based systems offload file storage to \ac{IPFS}. Applications typically write only the \ac{CID} and basic metadata on-chain while the image, document, or other payload lives in \ac{IPFS}. When the file is needed later, a user or service fetches it through its \ac{CID} and verifies integrity by recomputing the content hash. This split keeps ledger entries small and affordable while still giving tamper-evident links between on-chain records and off-chain data.

\section{Zero-Knowledge and Identity Concepts}

Zero-knowledge proofs allow a user to prove that they possess sensitive information without actually revealing the information itself. 

The identity decoupling is critical requirment because citizens may hesitate to report sensitive problems they have if their identity is linked to the submitted report. Therefore we decouple identity of the citizen by using public-private key pair associated with the citizen. By leveraging this approach the blockchain ledger will record only pseudonymous identifiers, authorization tickets and cryptographic proofs. This ensure no sensitive personal information is ever publicly exposed.

\section{AI-based Content Moderation}

Citizen reporting platforms are vulnerable to spam Submissions, abusive content, misinformation, and malicious reports. Due to immutability of blockchain records, permanently storing harmful content presents significant ethical and operational concerns. Therefore, content moderation techniques are required before reports are recorded on-chain. 

\ac{AI} techniques are widely used for automated content moderation. \ac{NLP} techniques are used in text moderation tasks such as:toxicity detection, hate speech detection, spam filtering, and sentiment Analysis. For image moderation, computer vision and deep learning approaches are mostly used. \ac{CNN}s and transformer-based vision models can classify images based on categories like: explicit content, graphic violence, hate symbols, and unrelated media. 

With these technical primitives established, Chapter 3~\ref{ch:architecture} details how they are integrated into a multi-layered architecture for local governance reporting system.
